\section{Berechnung der Zahnfußbeanspruchung}
\label{sec:berechnung_stand}

Die Zahnfußbeanspruchung eines Kunststoffrades wird nach dem Stand der Technik auf drei Wegen bestimmt. Mit der Richtlinie VDI 2736 Blatt 2 wird die Beanspruchung auf eine Nennspannung mit Einflussfaktoren zurückgeführt und mit einer temperatur- und lastspielzahlabhängigen Festigkeit verglichen \cite[S.~11 bis 15]{VDI2736-2}. Numerische Verfahren wie RIKOR, der FE-Baustein der FVA-Workbench oder eigene FE-Modelle lösen die Lastverteilung und den Spannungsverlauf in der Fußrundung auf und erlauben damit auch Fußgeometrien jenseits der Trochoide. Der vereinfachte Festigkeitsnachweis nach Stommel et al. bewertet die FE-Spannungen an jedem Knoten mit Kennwerten aus Kurzzeit-Spannungs-Dehnungs-Kurven \cite[Kap.~6]{Sto18}. Die drei Wege werden in dieser Reihenfolge dargestellt; die eigene Umsetzung steht in \secref{sec:analytik}, \secref{sec:werkzeugkette} und \secref{sec:auswertung_beanspruchung}.

\subsection{Berechnung nach VDI 2736}
\label{sec:vdi2736}

Ein rechnerischer Festigkeitsnachweis stellt der auftretenden Spannung eine ertragbare Spannung gegenüber. Der Quotient aus Bauteilfestigkeit und Bauteilbeanspruchung ist die Sicherheit, die eine Mindestsicherheit nicht unterschreiten darf \cite[S.~58]{Nie19}. Im Maschinenbau hat sich dafür das Nennspannungskonzept durchgesetzt, bei dem Baugröße, Kerbform und Kerbempfindlichkeit auf der Seite der Bauteilfestigkeit berücksichtigt werden \cite[S.~58]{Nie19}. Die Richtlinie VDI 2736 Blatt 2 überträgt dieses Konzept auf thermoplastische Stirnräder und folgt dabei dem Verfahren C der DIN 3990 mit Kraftangriff am Zahnkopf \cite[S.~36]{Has18}. Ausgangspunkt ist die Zahnfußspannung bei Kraftangriff am Zahnkopf unter der Annahme reiner Biegebeanspruchung, und der Nachweis ist für beide Räder zu führen \cite[S.~11]{VDI2736-2}. Mit der Nennumfangskraft $F_{\mathrm{t}}$ am Teilzylinder, der tragenden Zahnbreite $b$ und dem Normalmodul $m_{\mathrm{n}}$ gilt nach Gleichung (10) der Richtlinie \cite[S.~11]{VDI2736-2}
\begin{equation}
\sigma_{\mathrm{F}} = K_{\mathrm{F}}\,Y_{\mathrm{Fa}}\,Y_{\mathrm{Sa}}\,Y_{\varepsilon}\,Y_{\beta}\,\frac{F_{\mathrm{t}}}{b\,m_{\mathrm{n}}} \leq \sigma_{\mathrm{FP}} .
\label{eq:sigma_f_vdi}
\end{equation}
Der Faktor $K_{\mathrm{F}}$ fasst Anwendungsfaktor, Dynamikfaktor, Breitenfaktor und Stirnfaktor zusammen; wegen der Nachgiebigkeit der Thermoplaste gilt für $b/m_{\mathrm{n}} \leq 12$ vereinfacht $K_{\mathrm{F}} \approx K_{\mathrm{A}}$ \cite[S.~11]{VDI2736-2}. Den Einfluss der Zahnform auf die Biegenennspannung bei Lastangriff am Zahnkopf erfasst der Formfaktor $Y_{\mathrm{Fa}}$, die Spannungskonzentration durch die Kerbwirkung der Fußrundung der Spannungskorrekturfaktor $Y_{\mathrm{Sa}}$ \cite[S.~11 bis 12]{VDI2736-2}. Beide Faktoren stehen in Diagrammen über der Ersatzzähnezahl $z_{\mathrm{n}}$ für das Bezugsprofil nach DIN 867 mit $h_{\mathrm{fP}} = 1{,}25\,m_{\mathrm{n}}$ und $\rho_{\mathrm{fP}} = 0{,}25\,m_{\mathrm{n}}$ \cite[S.~12 bis 13]{VDI2736-2}. Für andere Profile verweist die Richtlinie auf die Berechnung nach DIN 3990-3 \cite[S.~12 bis 13]{VDI2736-2}. Der Überdeckungsfaktor $Y_{\varepsilon} = 0{,}25 + 0{,}75/\varepsilon_{\alpha}$ nach Gleichung (11) rechnet den Kopfangriff auf den geänderten Hebelarm und die geänderte Kerbwirkung des maßgebenden Kraftangriffs um, wobei die Profilüberdeckung ohne Verformung bestimmt wird \cite[S.~13]{VDI2736-2}. Mit dem Schrägenfaktor $Y_{\beta}$ nach Gleichung (12) gehen die günstigeren Eingriffsverhältnisse der Schrägverzahnung ein \cite[S.~13]{VDI2736-2}. Die zulässige Zahnfußspannung folgt nach Gleichung (13) der Richtlinie aus der Zahnfußgrenzfestigkeit $\sigma_{\mathrm{FG}}$ und der geforderten Mindestsicherheit \cite[S.~13]{VDI2736-2}:
\begin{equation}
\sigma_{\mathrm{FP}} = \frac{\sigma_{\mathrm{FG}}}{S_{\mathrm{Fmin}}} \qquad \text{mit} \qquad \sigma_{\mathrm{FG}} = Y_{\mathrm{St}}\,\sigma_{\mathrm{FlimN}} , \quad Y_{\mathrm{St}} \approx 2{,}0 .
\label{eq:sigma_fp_vdi}
\end{equation}

Ausgedrückt ist die Zeitschwellfestigkeit $\sigma_{\mathrm{FlimN}}$ als Nennspannung; sie hängt von der Zahnfußtemperatur $\vartheta_{\mathrm{F}}$ und der geforderten Lastwechselzahl $N_{\mathrm{L}}$ ab \cite[S.~14]{VDI2736-2}. Blatt 2 enthält sie als Diagramm und Tabelle für PA66 und POM im Bereich $N_{\mathrm{L}} = 10^{5}$ bis $10^{8}$, hinterlegt mit je einer Zahlenwertgleichung in Temperatur und Lastwechselzahl \cite[S.~14 bis 15]{VDI2736-2}. Für PA46 und für faserverstärkte Werkstoffe nennt die Richtlinie keine Werte. Die Zahnfußtemperatur wird nach Gleichung (7) der Richtlinie aus Umgebungstemperatur, Nennleistung, Reibwert, Zahnverlustgrad, Wärmeübergangsbeiwert und relativer Einschaltdauer geschätzt; der Ansatz geht auf Hachmann und Strickle zurück \cite[S.~8 und 30]{VDI2736-2}. Als Mindestsicherheit nennt die Richtlinie $S_{\mathrm{Fmin}} = 2{,}0$ für den Dauerbetrieb, zu rechnen mit $N_{\mathrm{L}} = 10^{8}$, und $S_{\mathrm{Fmin}} = 1{,}6$ für den zeitweisen Betrieb mit der tatsächlichen Lastwechselzahl \cite[S.~14]{VDI2736-2}. Der Wert ist nach Schadensrisiko und Genauigkeit der Annahmen festzulegen und zwischen Hersteller und Anwender zu vereinbaren \cite[S.~14]{VDI2736-2}. Zur Ausfallwahrscheinlichkeit der Festigkeitswerte macht Blatt 2 keine Angabe; Blatt 4 beschreibt die Umrechnung von Versuchsergebnissen auf eine Ausfallwahrscheinlichkeit von $10\,\%$ \cite[S.~22]{Has18}.

Hinter den Faktoren steht das Modell des am Fuß eingespannten Zahns. Der Dauerbruch geht in der Regel von der Zugseite nahe der Berührstelle der $30$°-Tangente aus \cite[S.~147]{NieWin03}. Deshalb wird die maximale örtliche Zugspannung dort als Maß für die Bruchgefährdung angesetzt und als Produkt aus Biegenennspannung und Spannungskorrekturfaktor berechnet \cite[S.~147]{NieWin03}. Linke bezeichnet denselben Faktor als Spannungskonzentrationsfaktor, der neben der Kerbwirkung den Schubanteil des Spannungszustands erfasst, und legt die Biegenennspannung an der $30$°-Tangente als Bezugsspannung fest \cite[S.~298 und 312]{Lin10}. ISO 6336-3 Methode B bezieht die Faktoren auf den Kraftangriff im äußeren Einzeleingriffspunkt. Die Sehne zwischen den Berührpunkten der $30$°-Tangenten an die Fußrundungen ist der maßgebende Querschnitt \cite[S.~14]{ISO6336-3}. Gültig sind die Formeln für Bezugsprofile mit $\rho_{\mathrm{fP}} > 0$, bei denen der Berührpunkt auf der vom Werkzeug erzeugten Fußrundung liegt \cite[S.~14]{ISO6336-3}. Der Tangentenwinkel $\theta$ wird iterativ aus Gleichung (29) der Norm bestimmt \cite[S.~19]{ISO6336-3}. Daraus folgen die Zahnfußdickensehne $s_{\mathrm{Fn}}$, der Krümmungsradius $\rho_{\mathrm{F}}$ an der $30$°-Tangente und der Biegehebelarm $h_{\mathrm{Fe}}$ nach den Gleichungen (30) bis (32) \cite[S.~19]{ISO6336-3}. Mit dem Kraftangriffswinkel $\alpha_{\mathrm{Fen}}$ lautet der Formfaktor nach Gleichung (9) der Norm \cite[S.~17]{ISO6336-3}
\begin{equation}
Y_{\mathrm{F}} = \frac{6\,\dfrac{h_{\mathrm{Fe}}}{m_{\mathrm{n}}}\cos\alpha_{\mathrm{Fen}}}{\left(\dfrac{s_{\mathrm{Fn}}}{m_{\mathrm{n}}}\right)^{2}\cos\alpha_{\mathrm{n}}}\,f_{\varepsilon}
\label{eq:y_f_iso}
\end{equation}
wobei $f_{\varepsilon} = 1$ für Geradverzahnungen mit einer Profilüberdeckung unter $2$ gilt. Der Spannungskorrekturfaktor folgt nach den Gleichungen (62) bis (64) der Norm aus dem Kerbparameter $q_{\mathrm{s}}$ und dem Verhältnis $L$ und gilt für $1 \leq q_{\mathrm{s}} < 8$ \cite[S.~24 bis 25]{ISO6336-3}:
\begin{equation}
Y_{\mathrm{S}} = \left(1{,}2 + 0{,}13\,L\right) q_{\mathrm{s}}^{\,1/(1{,}21 + 2{,}3/L)} \qquad \text{mit} \qquad L = \frac{s_{\mathrm{Fn}}}{h_{\mathrm{Fe}}} , \quad q_{\mathrm{s}} = \frac{s_{\mathrm{Fn}}}{2\,\rho_{\mathrm{F}}} .
\label{eq:y_s_iso}
\end{equation}
Für den Kraftangriff am Zahnkopf werden dieselben Beziehungen mit dem Hebelarm des Kopfangriffs ausgewertet und liefern $Y_{\mathrm{Fa}}$ und $Y_{\mathrm{Sa}}$; die Umrechnung auf den äußeren Einzeleingriffspunkt übernimmt $Y_{\varepsilon}$ \cite[S.~148 bis 149]{NieWin03}.

Die Faktoren setzen starre Zähne voraus. Bei Kunststoffrädern verformen sich die Zähne unter Last so stark, dass das folgende Zahnpaar vorzeitig in den Eingriff kommt und das vorhergehende nachzeitig austritt. Die unter Last wirksame Profilüberdeckung $\varepsilon_{\alpha,\mathrm{w}}$ liegt über der theoretischen, und der Übergang in das Einzeleingriffsgebiet erfolgt später bei kleinerem Hebelarm \cite[S.~31]{Has18}. Hasl zeigt mit linear-elastischen FE-Rechnungen an der Paarung Stahl-Kunststoff, dass die Zahnfußspannung nach VDI 2736 bei kleiner Last gut mit der FE übereinstimmt \cite[S.~29 bis 31]{Has18}. Mit steigender Last und Zähnezahl liegt die Richtlinie zunehmend darüber \cite[S.~31 bis 32]{Has18}. Sein Verfahren ACORA bestimmt $\varepsilon_{\alpha,\mathrm{w}}$ iterativ aus der Eingriffsfedersteifigkeit und der hinreichenden Verformung in Eingriffsrichtung nach Thoma \cite[S.~35 bis 39]{Has18}. Mit einem daraus gebildeten Überdeckungsfaktor liegt die Zahnfußnennspannung nahe an FE und RIKOR \cite[S.~43 bis 44]{Has18}. Weil die Zeitschwellfestigkeitswerte der Richtlinie mit dem Ansatz starrer Zähne aus Laufversuchen zurückgerechnet wurden, enthalten sie den Überdeckungseinfluss der Prüfräder \cite[S.~47]{Has18}. Hasl berücksichtigt ihn beim Übergang auf das modifizierte Verfahren nachträglich über einen Überdeckungsfaktor der Prüfräder \cite[S.~48]{Has18}. Fürstenberger hatte zuvor einen erweiterten Überdeckungsfaktor mit $\varepsilon_{\alpha,\mathrm{w}}$ und einem gemittelten Korrekturfaktor von etwa $1{,}2$ vorgeschlagen \cite[S.~132 bis 133]{Fuer13}.

\subsection{Numerische Verfahren und Zahnfußgeometrien}
\label{sec:numerik_fussgeometrien}

Das FVA-Programm RIKOR wurde zur Auslegung von Flankenkorrekturen entwickelt. Es berechnet die Lastverteilung im Zahneingriff aus den elastischen Verformungen des Getriebeumfelds, der Zähne durch Biegung, Schub und Hertzsche Abplattung sowie aus Wellenverlagerungen, Fertigungsabweichungen und vorhandenen Korrekturen \cite[S.~1]{FVA30-XI}. Die ebene Betrachtung mittelt die Zahnsteifigkeit über den Eingriff und ist in das Verfahren für den Breitenfaktor der DIN 3990 eingegangen \cite[S.~1]{FVA30-XI}. Daneben liefert eine topologische Betrachtung aus den örtlichen Steifigkeiten der Flankenbereiche die Lastverteilung über den gesamten Eingriff; darauf stützen sich die hinterlegten Verfahren für Zahnfußbruch, Grübchen und Graufleckigkeit \cite[S.~1 bis 3]{FVA30-XI}. Für die Zahnfußspannungsverteilung dient die Biegemomentverteilung einer einseitig eingespannten Platte nach Umezawa, eingespannt am Berührpunkt der $30$°-Tangente \cite[S.~43]{FVA30-XI}. Mit dem Spannungskorrekturfaktor nach DIN 3990 aus Zahnfußdickensehne, Fußrundungsradius und örtlichem Hebelarm wird die Biegenennspannung über der Zahnbreite zur örtlichen Spannung umgerechnet \cite[S.~43 bis 44]{FVA30-XI}. Das aus FVA 732 stammende Modul der örtlichen Zahnfußspannung ergänzt die Kerbspannungsverläufe für alle Punkte der Fußkurve auf Zug- und Druckseite und für Kraftangriff vom Kopf bis zum Fuß \cite[S.~11 bis 12]{FVA30-XI}. Damit lassen sich auch besondere Fußgeometrien bewerten \cite[S.~12]{FVA30-XI}. Seit der Version L wird das Gesamtsystem eingriffsstellungsabhängig durchgewälzt, mit vor- und nachzeitigem Eingriff \cite[S.~11 bis 12]{FVA30-XI}. Das Vorhaben FVA 30 XII ersetzt den analytischen Halbraumansatz nach Weber und Banaschek für den Radkörper durch ein FE-Modell \cite[S.~3]{FVA30-XII}. Dessen Arbeitsintegrale an der Schnittfläche zwischen Zahn und Radkörper gehen in die Eingriffssteifigkeit ein \cite[S.~3 bis 4]{FVA30-XII}. Für Kunststoffverzahnungen berücksichtigt RIKOR den Einfluss von Temperatur, Feuchte und Belastungsgeschwindigkeit auf das Werkstoffverhalten nicht \cite[S.~2]{FVA892-I}.

Die FVA-Workbench bündelt die Zahnkontaktanalysen STIRAK und RIKOR; ihre FE-basierte Analyse rechnete bis zum Vorhaben FVA 892 I nur linear-elastisch \cite[S.~1 bis 2]{FVA892-I}. Das Vorhaben erweitert die Workbench um eine quasistatische Abwälzsimulation auf Basis von ABAQUS \cite[S.~4]{FVA892-I}. Sie erfasst materielle und geometrische Nichtlinearitäten, den vor- und nachzeitigen Eingriff, die Kontaktfindung im verformten Zustand und einen beliebig geformten Radkörper \cite[S.~4]{FVA892-I}. Das Werkstoffverhalten wird über eine Spannungs-Dehnungs-Kurve als hyperelastisches Modell vorgegeben; ein Tragfähigkeitsnachweis und Wöhlerlinien sind ausdrücklich nicht Ziel des Vorhabens \cite[S.~5]{FVA892-I}. Ausgewertet werden an der Oberfläche jedes Zahnfußes aus dem Spannungsverlauf jedes Knotens über dem Abwälzen die Mittelspannung, die Ausschlagspannung und das Spannungsverhältnis \cite[S.~19]{FVA892-I}. Hinzu kommen eine schädigungsäquivalente Schwellspannung mit einer vom Nutzer vorgegebenen Mittelspannungsempfindlichkeit und die Dehnung im Fußbereich \cite[S.~19 und 21]{FVA892-I}. In einer Beispielrechnung mit der Spannungs-Dehnungs-Kurve des Kunststoffs Delrin 100 NC010 stimmen die Fußspannungen des Ritzels gut mit VDI 2736 und ISO 6336 überein \cite[S.~27 bis 28]{FVA892-I}. Die Fußspannungen des Rades weichen deutlich ab, was auf die Überdeckungsverhältnisse und den Radkörper zurückgeführt wird \cite[S.~29 bis 31]{FVA892-I}. Das Materialmodell beeinflusst Fußspannung, Dehnung, Gesamtüberdeckung und Rechenzeit entscheidend \cite[S.~31]{FVA892-I}. Der Radkörper verändert die Eingriffssteifigkeit; Querbohrungen können sie erheblich mindern, während die Normen den Radkörper nur über die Kranzdicke und RIKOR über einen Radkörperfaktor in Abhängigkeit der Stegbreite berücksichtigen \cite[S.~2 bis 3]{FVA484-V}. FVA 484 V bindet einen aus CAD importierten, mit Tetraedern vernetzten Radkörper an das Verzahnungsnetz von STIRAK an \cite[S.~33]{FVA484-V}. Die über dem Umfang veränderliche Steifigkeit verlangt eine Bewertung über eine volle Umdrehung \cite[S.~34]{FVA484-V}. Diese Anbindung nutzt auch FVA 892 I \cite[S.~4]{FVA892-I}.

Die FE-Berechnung der Zahnfußspannung bildet den Zahn mit seiner Fußrundung ab und wertet die Spannung an der Oberfläche der Fußrundung aus. Frühe bestätigt an neun Verzahnungen die gute Übereinstimmung von FE und RIKOR und dokumentiert die Abweichungen der Normen bei großen Schrägungswinkeln und Profilüberdeckungen \cite[S.~104]{Frue12}. Für Kunststoffpaarungen setzt Fürstenberger ein dreidimensionales ANSYS-Modell mit 60 Eingriffsstellungen ein und wertet die erste Hauptspannung aus \cite[S.~36 bis 37]{Fuer13}. Bei verlustoptimierten Verzahnungen berechnet VDI 2736 bis zu $30\,\%$ zu hohe Fußspannungen, weil der Schadensausgang bei größeren Tangentenwinkeln liegt und die Überdeckung unter Last fehlt \cite[S.~94 bis 95]{Fuer13}. Hasl rechnet linear-elastisch mit 20 Eingriffsstellungen und versteht unter der FE-Zahnfußspannung die über alle Stellungen maximale Hauptnormalspannung an der Zugseite der Fußrundung \cite[S.~29]{Has18}. Kassem et al. modellieren das Stahlrad als starre Schale, die halbe Zahnbreite mit Symmetriebedingung und PA46-GF30 mit einem orthotrop elastisch-plastischen Hill-Modell samt nachgebildeter Faserorientierung \cite[S.~6 bis 9]{Kas22}. Als Zahnfußspannung gilt dort die Biegespannung in einem lokalen, der Fußoberfläche folgenden Koordinatensystem \cite[S.~7]{Kas22}. Im Vergleich mit VDI 2736 stimmt die Trochoide bei PA46 für den Modul $1$\,mm gut überein, für Modul $2$\,mm überschätzt die Richtlinie die FE-Spannung um etwa $20\,\%$ \cite[S.~12]{Kas22}. Bei PA46-GF30 unterschätzt sie die FE-Spannung für Modul $1$\,mm um bis zu $25\,\%$ \cite[S.~14]{Kas22}. Analytische Verfahren stoßen nach Kassem et al. an ihre Grenzen, sobald inhomogene Werkstoffe eingehen \cite[S.~17]{Kas22}.

Die Spritzgießform erlaubt Fußrundungen, die nicht an ein Wälzwerkzeug gebunden sind. Frühe leitet die Gleichungen einer Zahnfußellipse her, deren erste Hauptachse um den Winkel $\gamma$ gegen die Zahnfußhalbierende geneigt ist \cite[S.~93 bis 94]{Frue12}. Sie geht am Fußformkreispunkt $P_{\mathrm{Ff}}$ in die Evolvente und am Fußpunkt $P_{\mathrm{f}}$ in die Lückenmitte tangentenstetig über \cite[S.~93]{Frue12}. Von den fünf Freiheitsgraden der Ellipse bleiben so die Neigung $\gamma$ und das Achsenverhältnis $a/b$ frei; Hauptachsen und Fußkreisdurchmesser folgen aus der Tangentenstetigkeit \cite[S.~94 bis 98]{Frue12} (\figref{fig:quelle_fruehe_ellipse}). In einer FE-Parameterstudie an 19 Referenzverzahnungen variiert Frühe $\gamma$ von $20$ bis $35$° und $a/b$ von $2{,}0$ bis $4{,}25$ \cite[S.~99]{Frue12}. Die Ellipse mit $\gamma = 30$° und $a/b = 3{,}0$ senkt die Fußspannung im Mittel um $11{,}9\,\%$ gegenüber dem bereits optimierten Standardwerkzeug, das Optimum aller Ellipsen um $14{,}5\,\%$ \cite[S.~100 bis 101]{Frue12}. Ellipsenähnliche Formen bringen keinen Zugewinn, und ein freier Spline im Optimierer liegt höchstens $3\,\%$ darunter \cite[S.~101]{Frue12}. Die analytische Rechnung an der $30$°-Tangente weicht für solche Fußformen um $12$ bis $30\,\%$ von der FE ab \cite[S.~103]{Frue12}. Wegen des größeren hoch beanspruchten Volumens überträgt sich die Spannungsminderung nicht im gleichen Verhältnis auf die Tragfähigkeit \cite[S.~102]{Frue12}. Kassem et al. legen die Ellipse über den Fußkreisdurchmesser fest und variieren diesen mit der Fußhöhe des Bezugsprofils \cite[S.~3 bis 4]{Kas22}. Bei PA46-GF30 liegt das Optimum bei kleinem Fußhöhenfaktor und einem um $0{,}2$ bis $0{,}4$ Modul verkleinerten Fußkreisdurchmesser \cite[S.~14]{Kas22}. Landi et al. passen Kreisbogen und Ellipse zwischen einem Punkt unterhalb des Nutzkreises und der Lückenmitte ein und werten Formfaktor und Spannungskorrekturfaktor punktweise entlang der Kontur aus \cite[S.~7 bis 10]{Lan21}. Gegenüber der Trochoide finden sie eine um etwa $15\,\%$ geringere Biegespannung, bestätigt durch FE mit weniger als $5\,\%$ Abweichung \cite[S.~11 und 16 bis 17]{Lan21}. Sanders bestätigt mit Dehnungsmessstreifen an Stahlrädern die mit dem Programm LDP vorhergesagte Minderung der Fußspannung um $22{,}4\,\%$ durch eine elliptische Fußrundung und um $24{,}3\,\%$ in Verbindung mit asymmetrischen Zähnen \cite[S.~iii]{San10}. In Pulsatorversuchen stieg die mittlere Lebensdauer dabei um mehr als das Dreißigfache \cite[S.~iv]{San10}.

\begin{figure}[htbp]
    \centering
    \includegraphics[width=0.5\textwidth]{kap2/quelle_fruehe_ellipse}
    \caption{Zahnfußausrundung mit tangentenstetigen Ellipsen, Neigung $\gamma$ und Hauptachsen $a$ und $b$, nach \cite[S.~93]{Frue12}.}
    \label{fig:quelle_fruehe_ellipse}
\end{figure}

Voith beschreibt in der Patentschrift DE 10 2013 004 861 B3 eine Fußrundung aus einer einzigen kubischen Bézierkurve \cite[S.~2]{DE102013004861B3}. Sie geht ab einem relevanten Durchmesser $d_{\mathrm{r}}$ in beiden Hauptpunkten tangentenstetig in das Zahnprofil über \cite[S.~3]{DE102013004861B3}. Weil die Krümmung stetig verläuft, entsteht kein Krümmungssprung, und die Spannungen verteilen sich gleichmäßiger \cite[S.~3]{DE102013004861B3}. Die beiden Kontrollpunkte liegen auf den Tangenten in den Hauptpunkten \cite[S.~4]{DE102013004861B3} (\figref{fig:quelle_patent_bezier}). Ihr Abstand vom Hauptpunkt beträgt $(0{,}25 + 0{,}1\,f)\,l$ mit dem Abstand $l$ des Hauptpunkts vom Schnittpunkt der Tangenten und dem Faktor $0 \leq f \leq 3$, bevorzugt $0{,}5$ bis $1{,}5$ \cite[S.~4]{DE102013004861B3}. Bei symmetrischer Lücke bleibt damit ein einziger Parameter \cite[S.~3]{DE102013004861B3}. Für ein Beispielrad mit Modul $4$\,mm und 50 Zähnen nennt die Schrift eine Verbesserung von rund $35\,\%$ gegenüber der Verzahnung nach DIN 867 \cite[S.~8 bis 9]{DE102013004861B3}. Roth und Opferkuch beschreiben dieselbe Kontur mit einem Bézier-Faktor, der die Lage der Kontrollpunkte auf der Kathete des ersten Zugdreiecks nach Mattheck festlegt \cite[S.~722]{RotOpf17}. An einem Ritzel mit 25 Zähnen und Modul $6{,}5$\,mm liegt das Optimum bei $0{,}57$ mit einer theoretischen Steigerung der Zahnfußsicherheit um $24\,\%$, im Bereich $0{,}46$ bis $0{,}81$ über $20\,\%$ \cite[S.~724 bis 725]{RotOpf17}. Über ein Feld von Zähnezahlen und Profilverschiebungen ergeben sich Steigerungen von mindestens $12\,\%$ bis zu $30\,\%$, wobei die Bézierkontur die tragende Flanke unberührt lässt \cite[S.~728 bis 729]{RotOpf17}. Im Vorhaben FVA 241 XII wurde eine so optimierte Prüfverzahnung mit Modul $4$\,mm gefertigt. Die FE-Studie zeigt für Bézier-Faktoren von $0{,}5$ bis $0{,}6$ eine Minderung der örtlichen Fußspannung um etwa $25\,\%$ \cite[S.~43]{FVA241-XII}. Im Pulsatorversuch trägt die Variante mit dem Faktor $0{,}57$ eine um etwa $20\,\%$ höhere dauerfest ertragbare Normalkraft \cite[S.~93]{FVA241-XII}. Die Zahnfußdauerfestigkeiten nach ISO 6336-3 und nach RIKOR weichen für diese Kontur um höchstens $3\,\%$ voneinander ab, mit dem Spannungsmaximum bei einem Tangentenwinkel von $36$ bis $37$° \cite[S.~94 bis 95]{FVA241-XII}. Dong et al. verfolgen den umgekehrten Weg und beschreiben die Kopfkontur des Wälzfräsers mit einer Bézierkurve aus fünf Punkten \cite[S.~4]{Don20}. Die Optimierung mit genetischem Algorithmus und FE senkt die maximale Fußspannung dreier Stahlräder um $18{,}19$, $20{,}06$ und $22{,}02\,\%$ \cite[S.~12 und 15]{Don20}.

\begin{figure}[htbp]
    \centering
    \includegraphics[width=0.5\textwidth]{kap2/quelle_patent_bezier}
    \caption{Fußrundung aus einer kubischen Bézierkurve mit den Hauptpunkten $P_{0}$ und $P_{3}$, den Kontrollpunkten $P_{1}$ und $P_{2}$ auf den Tangenten und dem Schnittpunkt $S$, nach \cite[Fig.~2]{DE102013004861B3}.}
    \label{fig:quelle_patent_bezier}
\end{figure}

Auf die Zugdreiecke nach Mattheck geht die bionische Fußrundung zurück, die in der Patentschrift DE 10 2008 045 318 B3 beschrieben ist. Im Koordinatensystem der Zahnlücke mit der Symmetrieachse als Ordinate geht die Evolvente am relevanten Durchmesser $d_{\mathrm{r}}$ tangentenstetig in eine Tangensfunktion über \cite[S.~3 bis 4]{DE102008045318B3}. Diese mündet tangentenstetig in einen Kreisbogen, der im Fußpunkt am Fußkreis anliegt \cite[S.~4]{DE102008045318B3} (\figref{fig:quelle_patent_bionisch}). Der relevante Durchmesser wird als arithmetisches Mittel aus Nutzkreis $d_{\mathrm{N}}$ und Freikreis $d_{\mathrm{FR}}$ gewählt, damit der Gegenzahn stets auf der Evolvente abläuft \cite[S.~3]{DE102008045318B3}. Frei bleiben der Steigungswinkel $\varphi$ der Tangensfunktion auf der Symmetrieachse, im Patent $\gamma$ genannt, und der Radius $r$ des Kreisbogens \cite[S.~4 und 6]{DE102008045318B3}. Der Winkel soll unter $65$° liegen, das Optimum nahe oder unter $45$°, bevorzugt nach der Beziehung $45^{\circ} - 180^{\circ}/z - \alpha_{\mathrm{r}}$ mit dem Eingriffswinkel $\alpha_{\mathrm{r}}$ am relevanten Durchmesser \cite[S.~5 und 9]{DE102008045318B3}. Für den Radius haben sich das $0{,}1$- bis $0{,}6$-Fache, bevorzugt das $0{,}3$- bis $0{,}4$-Fache der Lückenweite $S_{\mathrm{L}}$ am relevanten Durchmesser bewährt \cite[S.~7]{DE102008045318B3}. FE-Rechnungen der Erfinder ergaben Steigerungen der Zahnfußfestigkeit von bis zu $30\,\%$ gegenüber dem wälzgefrästen Rad \cite[S.~2 und 7]{DE102008045318B3}. Kassem et al. übernehmen diese Kontur mit dem Radius als Vielfachem von $S_{\mathrm{L}}$ \cite[S.~5 bis 6]{Kas22}. Bei PA46-GF30 liegt das Optimum beider Baugrößen bei einem Winkel von etwa $30$° und kleinem Radius \cite[S.~14]{Kas22}. Für Modul $2$\,mm führt die bionische Fußrundung zur niedrigsten Spannung; die Trochoide liegt um $7{,}3\,\%$, die Ellipse um $4{,}84\,\%$ und die Vollrundung um $2{,}63\,\%$ darüber \cite[S.~16]{Kas22}. Die optimierten Konturen unterscheiden sich untereinander nur wenig, verstärkter und unverstärkter Werkstoff verlangen aber unterschiedliche Parameter \cite[S.~16 bis 17]{Kas22}. Im Verschleiß verlagert sich das Maximum in die Verschleißkerbe, in der Trochoide, Ellipse und bionische Kontur zusammenfallen \cite[S.~16 bis 17]{Kas22}. In einer Folgearbeit lassen Kassem et al. die Fußkontur nach der Wachstumsregel der biologischen Formoptimierung CAO wachsen \cite[S.~2596]{Kas23}. Knoten mit Spannung über einer Referenzspannung werden nach außen, darunter nach innen verschoben, mit höchstens $0{,}025\,m_{\mathrm{n}}$ je Schritt und Kontrolle des Gegenzahns über den gesamten Eingriff \cite[S.~2596 bis 2597]{Kas23}. Gegenüber der Trochoide mit größtmöglicher Ausrundung sinkt die Fußspannung um $17$ bis $24\,\%$, wobei Werkstoff, Reibwert und gewählte Vergleichsspannung die optimale Kontur verändern \cite[S.~2598 bis 2600]{Kas23}.

\begin{figure}[htbp]
    \centering
    \includegraphics[width=0.5\textwidth]{kap2/quelle_patent_bionisch}
    \caption{Bionische Fußrundung aus Tangensfunktion und Kreisbahn ab dem relevanten Durchmesser $d_{\mathrm{r}}$ mit dem Steigungswinkel auf der Symmetrieachse und dem Übergangspunkt zur Kreisbahn, nach \cite[Fig.~2]{DE102008045318B3}.}
    \label{fig:quelle_patent_bionisch}
\end{figure}

\subsection{Bemessung nach Stommel et al.}
\label{sec:stommel_stand}

Stommel, Stojek und Korte schlagen einen vereinfachten Festigkeitsnachweis für thermoplastische Kunststoffbauteile vor, der ausschließlich auf kurzzeitigen isothermen und langzeitigen isochronen Spannungs-Dehnungs-Kurven aufsetzt und auf den Bemessungsansatz nach Oberbach zurückgeht \cite[S.~308]{Sto18}. Das Verfahren ist phänomenologisch, hat abschätzenden Charakter und ist für Machbarkeitsstudien und frühe Entwicklungsphasen gedacht \cite[S.~307 bis 309]{Sto18}. Die Autoren haben den Ansatz um die Festlegung eines eindeutigen Abminderungsfaktors und um einen weitgehend neuen Ermüdungsnachweis erweitert und in eine formelle Struktur bis zum Auslastungsgrad gebracht \cite[S.~309]{Sto18}. Grundlage ist das Konzept der lokalen Bauteilfestigkeit. Spannungsgradienten, Mehraxialität, Bindenähte und Orientierungen senken die Festigkeit örtlich verschieden, so dass die Stelle der höchsten Beanspruchung nicht die Stelle des höchsten Auslastungsgrads sein muss \cite[S.~310 bis 313]{Sto18}. Der Auslastungsgrad ist der Quotient aus vorliegender und zulässiger Beanspruchung, der Kehrwert des Sicherheitsfaktors, und kann spannungs-, dehnungs- oder energiebezogen gebildet werden \cite[S.~311 und 313]{Sto18}. Die Auswahl der Versagenshypothese steuert, formell analog zur FKM-Richtlinie, eine Interaktionsgleichung zwischen Normalspannungshypothese und von-Mises-Spannung \cite[S.~322]{Sto18}.

Aus dem Schubfestigkeitsfaktor $f_{\tau}$, dem Verhältnis von Schub- zu Zugfestigkeit, folgt die Steuergröße $q$ dieser Interaktionsgleichung nach Gleichung (6.2) des Buches \cite[S.~322]{Sto18}
\begin{equation}
q = \frac{\sqrt{3} - 1/f_{\tau}}{\sqrt{3} - 1} .
\label{eq:q_ftau}
\end{equation}
Ideal zähes Verhalten entspricht $f_{\tau} = 0{,}5$ und ideal sprödes $1{,}0$; ein Versagen allein aus den gestaltändernden Anteilen ergibt $0{,}58$ \cite[S.~323]{Sto18}. Tabelle 6.1 des Buches nennt für Raumtemperatur und quasistatische Beanspruchung $f_{\tau} = 0{,}58$ und $q = 0$ für duktile teilkristalline Thermoplaste \cite[S.~323]{Sto18}. Für verstärkte teilkristalline Thermoplaste stehen dort $0{,}65$ und $0{,}26$, für amorphe $0{,}74$ und $0{,}52$ \cite[S.~323]{Sto18}. Feuchte hydrophile verstärkte Thermoplaste wie Polyamid mit Glasfaser ordnet die Fußnote der Tabelle der duktilen Spalte zu \cite[S.~323]{Sto18}. Den lokalen Beanspruchungszustand beschreibt der Mehraxialitätsgrad $h$ nach Hencky, Gleichung (6.3) des Buches \cite[S.~325]{Sto18}:
\begin{equation}
h = \frac{\sigma_{\mathrm{H}}}{\sigma_{\mathrm{V,Mises}}} = \frac{(\sigma_{1} + \sigma_{2} + \sigma_{3})/3}{\sqrt{\tfrac{1}{2}\left[(\sigma_{1} - \sigma_{2})^{2} + (\sigma_{2} - \sigma_{3})^{2} + (\sigma_{3} - \sigma_{1})^{2}\right]}} .
\label{eq:h_hencky}
\end{equation}
Reiner Schub ergibt $h = 0$, der einachsige Zug $1/3$ und der hydrostatische Zustand unendlich; an der kräftefreien Oberfläche ist $2/3$ der Höchstwert \cite[S.~325]{Sto18}. Mit steigendem $h$ versagt der Werkstoff spröder bei geringerer Dehnung \cite[S.~324 bis 325]{Sto18}. Die zulässige Dehnung wird deshalb nach Gleichung (6.4) des Buches zwischen einer oberen Grenze $\varepsilon_{\mathrm{max}}$ und einer unteren Grenze $\varepsilon_{\mathrm{min}}$ interpoliert \cite[S.~326]{Sto18}:
\begin{equation}
\varepsilon_{\mathrm{zul}} =
\begin{cases}
\varepsilon_{\mathrm{max}} & \text{für } h \leq 1/3 ,\\[4pt]
\left(1 - \dfrac{1}{3h}\right)\left(\varepsilon_{\mathrm{min}} - \varepsilon_{\mathrm{max}}\right) + \varepsilon_{\mathrm{max}} & \text{für } h > 1/3 .
\end{cases}
\label{eq:eps_zul}
\end{equation}
Beide Grenzen liefert das Nomogramm in \figref{fig:nomogramm_qs}. Es ist auf den versagenskritischen Werkstoffkennwert $K$ der Kurve bei Raumtemperatur normiert, die Streckgrenze oder bei fehlender Streckgrenze die Bruchspannung \cite[S.~344 bis 346]{Sto18}. Je Werkstoffklasse und Lastart ordnet es einen Bereich des Abminderungsfaktors $A$ zu \cite[S.~345 bis 346]{Sto18}. Für die einmalige kurzzeitige Last reicht der Bereich von $0{,}8$ bis $1{,}0$ für unverstärkte teilkristalline und von $0{,}55$ bis $0{,}70$ für verstärkte Thermoplaste \cite[S.~346]{Sto18}. Die kleineren Faktoren der verstärkten Klassen erfassen pauschal, dass die Spannungen isotrop und mit der günstigsten Kurve in Faserrichtung ermittelt wurden und dass die plastische Reserve fehlt \cite[S.~345 bis 346]{Sto18}. Obere und untere Grenze schneiden die Raumtemperaturkurve bei $\varepsilon_{\mathrm{max}}$ und $\varepsilon_{\mathrm{min}}$. Die zulässige Dehnung gilt als temperaturunabhängig und wird senkrecht auf die Isotherme der Betriebstemperatur übertragen, von der die statische Bauteilfestigkeit $\sigma_{\mathrm{stat}}$ abgelesen wird \cite[S.~347 bis 350]{Sto18}. Negative Hauptspannungen werden vor der Bildung der Vergleichsspannung nach Gleichung (6.17) des Buches durch den Druckfestigkeitsfaktor geteilt, weil für Kunststoffe eine um $30\,\%$ höhere Druck- als Zugfestigkeit angenommen wird \cite[S.~357]{Sto18}:
\begin{equation}
\bar{\sigma}_{i} = \frac{\sigma_{i}}{f_{\sigma,i}} \qquad \text{mit} \qquad f_{\sigma,i} = 1 \text{ für } \sigma_{i} \geq 0 , \quad f_{\sigma,i} = 1{,}3 \text{ für } \sigma_{i} < 0 .
\label{eq:druckkorrektur}
\end{equation}
Die Vergleichsspannung und der dehnungsbezogene statische Auslastungsgrad stehen als \eqref{eq:a_stat} in \secref{sec:auswertung_beanspruchung}. Abgelesen wird die Vergleichsdehnung nicht aus dem Verzerrungstensor, sondern zur Vergleichsspannung auf der Spannungs-Dehnungs-Kurve, und Sicherheitsfaktoren sind im Verfahren implizit enthalten \cite[S.~355 bis 358]{Sto18}.

\begin{figure}[htbp]
    \centering
    \includegraphics[width=0.85\textwidth]{kap2/nomogramm_qs}
    \caption{Nomogramm der statischen Bemessung nach Stommel et al. für PA46-GF30 mit der quasistatischen Werkstoffkarte, in Anlehnung an \cite{Sto18}.}
    \label{fig:nomogramm_qs}
\end{figure}

Weil Kunststoffe keine Dauerfestigkeit kennen, dient die Spannungsamplitude bei der Grenzlastspielzahl $N_{\mathrm{G}} = 10^{7}$ als Ersatzdauerfestigkeit; jenseits davon ist mit weiterem Abfall zu rechnen \cite[S.~328]{Sto18}. Der Ermüdungsnachweis gilt nur oberhalb von $10^{4}$ Lastspielen, setzt einen erfolgreichen statischen Nachweis voraus und ermittelt die Spannungen stets linear-elastisch \cite[S.~361 bis 362]{Sto18}. Die Bauteil-Wechselfestigkeit $\sigma_{\mathrm{W}} = A_{\mathrm{W}} K$ folgt aus dem zyklischen Abminderungsfaktor desselben Nomogramms, der für $10^{7}$ Lastspiele bei $R = -1$ definiert ist \cite[S.~362 bis 363]{Sto18}. Für kurzfaserverstärkte Klassen ist er ein fester Wert, $0{,}25$ bei teilkristalliner Matrix, für unverstärkte teilkristalline Thermoplaste ein Bereich von $0{,}20$ bis $0{,}30$ \cite[S.~362 bis 363]{Sto18}. Der Faktor lässt sich als Wechselfestigkeitsfaktor auffassen, wie er für Stähle mit $0{,}4$ bis $0{,}45$ der Zugfestigkeit bekannt ist \cite[S.~363]{Sto18}. Für ein POM und ein PA66-GF liegen berechnete und gemessene Wechselfestigkeit nahe beieinander \cite[S.~363 bis 364]{Sto18}. Die Betriebsfestigkeit nach \eqref{eq:betriebsfestigkeit} korrigiert $\sigma_{\mathrm{W}}$ mit dem Mittelspannungsfaktor und dem Betriebsfestigkeitsfaktor \cite[S.~360]{Sto18}. Das Haigh-Diagramm in \figref{fig:haigh_qs_80C} wird aus der Streckspannung $\sigma_{\mathrm{Y}}$, der statischen Bauteilfestigkeit und der Wechselfestigkeit konstruiert. Die Goodman-Gerade von $\sigma_{\mathrm{W}}$ zu $\sigma_{\mathrm{Y}}$ trägt die Mittelspannungsempfindlichkeit $M = \sigma_{\mathrm{W}}/\sigma_{\mathrm{Y}}$ als Betrag ihrer Steigung, und $\sigma_{\mathrm{stat}}$ begrenzt das Diagramm auf beiden Achsen \cite[S.~364 bis 365]{Sto18}. Gewählt wurde die Goodman-Korrektur als die konservativere gegenüber Gerber, normiert auf die Kurzzeitfestigkeit statt auf die Zeitstandfestigkeit \cite[S.~328 bis 329]{Sto18}. Der Mittelspannungsfaktor $K_{\mathrm{M}} = \sigma'_{\mathrm{W}}/\sigma_{\mathrm{W}}$ ist das Verhältnis der mittelspannungskorrigierten zur ursprünglichen Wechselfestigkeit und hängt davon ab, auf welchem Weg der Beanspruchungspunkt die Grenzlinie erreicht \cite[S.~366 bis 367]{Sto18}. Die vier Überlastungsfälle nach Gleichung (6.21) des Buches lauten \cite[S.~367 bis 368]{Sto18}
\begin{equation}
\begin{aligned}
K_{\mathrm{M}} &= 1 - M\,\frac{\sigma_{\mathrm{m}}}{\sigma_{\mathrm{W}}} \;\text{(F1)} , \qquad &
K_{\mathrm{M}} &= \frac{1}{1 + M\,\sigma_{\mathrm{m}}/\sigma_{\mathrm{a}}} \;\text{(F2)} , \\
K_{\mathrm{M}} &= \frac{1 - M\,\sigma_{\mathrm{min}}/\sigma_{\mathrm{W}}}{1 + M} \;\text{(F3)} , \qquad &
K_{\mathrm{M}} &= \frac{1 - M\,\sigma_{\mathrm{max}}/\sigma_{\mathrm{W}}}{1 - M} \;\text{(F4)}
\end{aligned}
\label{eq:k_m_ueberlastung}
\end{equation}
für F1 mit konstanter Mittelspannung, F2 mit konstantem Spannungsverhältnis, F3 mit konstanter Minimalspannung und F4 mit konstanter Maximalspannung; die Wahl obliegt dem Anwender \cite[S.~366 bis 367]{Sto18}.

\begin{figure}[htbp]
    \centering
    \includegraphics[width=0.85\textwidth]{kap2/haigh_qs_80C}
    \caption{Haigh-Diagramm der zyklischen Bemessung für PA46-GF30 bei $80$\,\textdegree C mit der quasistatischen Werkstoffkarte, in Anlehnung an \cite{Sto18}.}
    \label{fig:haigh_qs_80C}
\end{figure}

Zwischen $10^{4}$ Lastspielen und $N_{\mathrm{G}}$ wird die Wöhlerlinie synthetisch gebildet. Das Nomogramm wird logarithmisch umgezeichnet, so dass seine Achse mit der auf $K$ normierten Amplitudenachse eines doppeltlogarithmischen Wöhler-Diagramms zusammenfällt \cite[S.~369 bis 370]{Sto18}. Von $\sigma_{\mathrm{stat}}$ bei $N = 1$ bis zur Grenzspannungsamplitude bei $N_{\mathrm{G}}$ verläuft die Linie nach der Basquin-Gleichung geradlinig, und die Mittelspannungskorrektur verschiebt sie parallel nach unten \cite[S.~370]{Sto18}. Der Betriebsfestigkeitsfaktor und der Neigungsexponent in \eqref{eq:k_m_k_b} folgen daraus; für verstärkte teilkristalline Thermoplaste ergeben sich Neigungsexponenten von $15{,}65$ bis $20{,}44$, was im Bereich gemessener Werte von $11{,}9$ bis $25{,}0$ liegt \cite[S.~371 bis 372]{Sto18}. Überschreitet die Betriebsfestigkeit die statische Grenze, wird sie nach Gleichung (6.23) auf $\sigma_{\mathrm{stat}} - \sigma_{\mathrm{m}}$ begrenzt \cite[S.~372]{Sto18}. Der zyklische Auslastungsgrad wird komponentenweise aus Einzelauslastungsgraden gebildet, weil sich im mehraxialen zyklischen Fall keine Vergleichsspannung bilden lässt; die Verknüpfung nach \eqref{eq:a_bv} nutzt die Steuergröße ohne Druckkorrektur \cite[S.~373 bis 374]{Sto18}. Stützzahlen bleiben unberücksichtigt, und als gut gilt eine Vorhersage der ertragbaren Lastspielzahl innerhalb eines Faktors von zwei bis drei \cite[S.~375]{Sto18}.

Für die Spannungsermittlung schreibt das Buch ein isotropes Materialmodell vor, auch wenn der Werkstoff kurzfaserverstärkt ist. Kalibriert wird es an der Kurzzeitkurve spritzgegossener Prüfstäbe, deren Fließrichtung in Stablängsrichtung liegt, womit für verstärkte Werkstoffe die vorteilhaftesten Kennwerte angenommen werden \cite[S.~343 bis 344]{Sto18}. Die Festigkeitsminderung quer zur Fließrichtung soll das Nachweisverfahren selbst über den niedrigen Abminderungsfaktor der faserverstärkten Klasse erfassen \cite[S.~344 bis 346]{Sto18}. Unzulässig für das Verfahren ist die explizite Berücksichtigung einer Faserorientierung mit einem anisotropen Modell, ebenso wie pauschale Reduktionsfaktoren auf Steifigkeit oder Kurve \cite[S.~344]{Sto18}. Entsprechend hängt der zyklische Abminderungsfaktor kurzfaserverstärkter Werkstoffe nicht vom Mehraxialitätsgrad ab \cite[S.~363]{Sto18}. Nach dem Buch ist das vereinfachte Verfahren die Alternative zum umfassenden Ermüdungsnachweis kurzfaserverstärkter Kunststoffe, der die Richtungsabhängigkeit auflöst und nur mit Software wie S-Life Plastics möglich ist \cite[S.~235 und 375]{Sto18}.

Die Grundlagen dieses umfassenden Nachweises stehen im dritten Kapitel des Buches. Bei kurzfaserverstärkten Thermoplasten hängt die Festigkeit zusätzlich von der Beanspruchungsrichtung relativ zur Faserhauptrichtung und vom lokalen Orientierungsgrad ab \cite[S.~219]{Sto18}. In Faserrichtung treten Faserenthaftung, Faserbruch und Matrixbruch auf, quer dazu fast immer Matrixbruch \cite[S.~221]{Sto18}. Richtungsabhängige Kennwerte werden an Proben aus filmangespritzten Platten unter $0$°, $90$° und $30$° oder $45$° zur Fließrichtung ermittelt und gelten je Orientierungsgrad \cite[S.~226 bis 227]{Sto18}. Für die Bewertung mehraxialer Zustände stellt das Buch das Maximalspannungskriterium, das Tsai-Hill-Kriterium und das Tsai-Wu-Kriterium gegenüber \cite[S.~229 bis 234]{Sto18}. Das Tsai-Hill-Kriterium ist ein pauschales Kriterium mit ellipsoider Versagensfläche, das die Wechselwirkung der Komponenten erfasst, Zug- und Druckfestigkeit aber nicht unterscheidet \cite[S.~231 bis 232]{Sto18}. In der dreidimensionalen Form nach Gleichung (3.36) des Buches lautet es mit den Zugfestigkeiten $X$, $Y$ und $Z$ in den Orthotropierichtungen und den Schubfestigkeiten $S_{ij}$ \cite[S.~233]{Sto18}
\begin{equation}
F\left(\sigma_{22} - \sigma_{33}\right)^{2} + G\left(\sigma_{33} - \sigma_{11}\right)^{2} + H\left(\sigma_{11} - \sigma_{22}\right)^{2} + 2\left(L\,\sigma_{12}^{2} + M\,\sigma_{13}^{2} + N\,\sigma_{23}^{2}\right) = 1
\label{eq:tsai_hill}
\end{equation}
mit $2F = 1/Y^{2} + 1/Z^{2} - 1/X^{2}$, $2G = 1/Z^{2} + 1/X^{2} - 1/Y^{2}$, $2H = 1/X^{2} + 1/Y^{2} - 1/Z^{2}$ sowie $2L = 1/S_{12}^{2}$, $2M = 1/S_{13}^{2}$ und $2N = 1/S_{23}^{2}$. Bei transversaler Isotropie verringert sich die Zahl der Festigkeiten auf vier, und die Wurzel der linken Seite ist der Auslastungsgrad \cite[S.~233]{Sto18}. Die quadratische Form entspricht dem Hill-Potential in \eqref{eq:hill}; die Zuordnung der Konstanten $L$ und $N$ zu den Schubkomponenten ist im Buch gegenüber der Schreibweise von Abaqus vertauscht. Zyklisch nutzt das Buch das Master-Wöhlerkurven-Konzept, nach dem das Verhältnis von zyklischer zu statischer Festigkeit näherungsweise unabhängig von der Faserrichtung ist \cite[S.~236 bis 237]{Sto18}. Daneben werden Wöhlerkurven über die Eigenwerte des Orientierungstensors interpoliert \cite[S.~238 bis 239]{Sto18}. Im lokalen Wöhlerkurven-Konzept werden die Faserorientierungen aus der Spritzgießsimulation mit einer Kopplungssoftware wie Converse elementweise an die FE übergeben \cite[S.~241 bis 242]{Sto18}. Die anisotropen Spannungen werden komponentenweise in den Faserhauptrichtungen mit lokalen Wöhlerkurven verglichen und die Einzelauslastungsgrade nach einer anisotropen Hypothese wie Tsai-Hill überlagert \cite[S.~243]{Sto18}. Dafür genügen im Bereich hoher Lastspielzahlen linear-elastisch ermittelte anisotrope Spannungen \cite[S.~242]{Sto18}. Auf diese Abschnitte stützt sich die eigene anisotrope Erweiterung des vereinfachten Nachweises in \secref{sec:auswertung_beanspruchung}, die die Hill-Vergleichsspannung \eqref{eq:hill} im Faserkoordinatensystem an die Stelle der von-Mises-Spannung setzt.

% Herkunft:
% Seitenzuordnung der Textauszüge (Druckseite): VDI2736-2 = PDF-Seite; ISO6336-3 = PDF − 8; Sto18 = PDF − 18;
%   Frue12 = PDF − 12; Has18 = PDF − 18; Fuer13 = PDF − 14 (Entwurf 01.03.2013); FVA30-XI = PDF − 10;
%   FVA892-I = PDF − 8; FVA484-V = PDF − 6; FVA241-XII = PDF − 10; NieWin03 = PDF − 13; Lin10 = PDF − 1;
%   Nie19 = PDF − 15; Kas22 nach der Fußzeile "22FTM13 n"; Kas23 Zeitschriftenseiten 2592 bis 2601; Lan21 und
%   Don20 artikelinterne Seiten; San10 Seiten der Arbeit (iii, iv); Patente: Seite n der Schrift (n/12, n/15);
%   RotOpf17: PDF-Seite n als Druckseite 718 + n angenommen (14 Seiten = 719 bis 732).
% Einleitung: VDI2736-2 S. 11 bis 15; Sto18 Kap. 6 (S. 307 bis 375).
% 2.5.1 Absatz 1: Nie19 S. 58 (Gl. 3.58 Sicherheit, Nennspannungskonzept); Has18 S. 36 (Methode C wie VDI 2736);
%   VDI2736-2 S. 11 (Gl. 10 = eq:sigma_f_vdi, Kopfangriff, beide Räder, K_F = K_A K_v K_Fβ K_Fα, K_F ≈ K_A);
%   S. 11 bis 13 (Y_Fa Bild 1 S. 12, Y_Sa Bild 2 S. 13, Fußnoten 5 und 6 DIN 3990-3, Gl. 11 und 12 S. 13,
%   Gl. 13 S. 13 = eq:sigma_fp_vdi, Y_St ≈ 2,0).
% Absatz 2: VDI2736-2 S. 14 (σ_FlimN, Bild 3 PA66/POM, S_Fmin 2,0 und 1,6, Vereinbarung), S. 14 Fußnote 8
%   (Zahlenwertgleichungen), S. 15 (Tabelle 5), S. 8 (Gl. 7 Zahnfußtemperatur), S. 30 (Anhang B Hachmann/Strickle);
%   PA46 im Text nicht enthalten (Suche in VDI2736-2.txt ohne Treffer); Has18 S. 22 (Ausfallwahrscheinlichkeit).
% Absatz 3: NieWin03 S. 147 (21.7.3.1, 30°-Tangente, Y_S erfasst Kerbe und Spannungszustand), S. 148 bis 149
%   (vereinfachter Ansatz Kopfangriff, Y_Fa, Y_Sa, Y_ε, Gl. 21.7/29 bis 32); Lin10 S. 298 (6.5.2.1 Biegenennspannung
%   an der 30°-Tangente als Bezugsspannung, Spannungskonzentrationsfaktor mit Schubanteil), S. 312 (Grundbeziehungen,
%   Formfaktor); ISO6336-3 S. 14 (6.1, Methode B, 30°-Tangenten, Bedingungen a bis e), S. 17 (Gl. 9 = eq:y_f_iso,
%   f_ε Gl. 10), S. 18 bis 19 (Gl. 26 bis 32), S. 24 bis 25 (Gl. 62 bis 64 = eq:y_s_iso, 1 ≤ q_s < 8).
% Absatz 4: Has18 S. 31 (ε_α,w, späterer Übergang, kleinerer Hebelarm), S. 29 bis 32 (FE gegen VDI, Bild 4 S. 30,
%   Bild 5 S. 33), S. 35 bis 39 (ACORA, Thoma, Bild 6 S. 39), S. 43 bis 44 (Y_ε,ACR Gl. 44, Vergleich mit FE/RIKOR,
%   Bild 10 S. 46), S. 47 bis 48 (Zeitschwellfestigkeit enthält Überdeckung der Prüfräder, Y_εT Gl. 46);
%   Fuer13 S. 132 bis 133 (Y_ε' Gl. 10.23 bis 10.25, f_Yε,mit ≈ 1,2).
% 2.5.2 Absatz 1 (RIKOR): FVA30-XI S. 1 (1.1 Zweck, Einflüsse, ebene und topologische Betrachtung), S. 2 (1.2),
%   S. 3 (1.4 Tragfähigkeitsverfahren), S. 11 bis 12 (FVA 732/I und II OeZS, FVA 30/X Durchwälzen, Version L),
%   S. 43 bis 44 (2.13 Umezawa, 30°-Tangente, Gl. 2.13.1 bis 2.13.3); FVA30-XII S. 3 (Problemstellung,
%   Weber/Banaschek, Halbraum), S. 3 bis 4 (Forschungsziel, Arbeitsintegrale); FVA892-I S. 2 (RIKOR, FVA 887 I).
% Absatz 2 (Workbench, Radkörper): FVA892-I S. 1 bis 2 (Ausgangssituation, STIRAK/RIKOR), S. 4 (Forschungsziel,
%   484 V), S. 5 (keine Wöhlerlinien, kein Tragfähigkeitsnachweis), S. 19 (3.3.4.1 Formel 3-3, Mittelspannungs-
%   empfindlichkeit vom Nutzer), S. 21 (Dehnung), S. 27 bis 31 (3.3.6 Beispiel Delrin 100 NC010, Zusammenfassung);
%   FVA484-V S. 2 (Bild 1 Querbohrungen, Kranzdicke DIN 3990-3/ISO 6336-3), S. 3 (Stand der Forschung,
%   Radkörperfaktor C_R), S. 33 bis 34 (Zusammenfassung, Bewertung über eine Umdrehung).
% Absatz 3 (FE gegen Norm): Frue12 S. 104 (Zusammenfassung FE/RIKOR/Normen); Fuer13 S. 36 bis 37 (3.5 RIKOR I,
%   ANSYS, 60 Eingriffsstellungen, 1. Hauptspannung), S. 94 bis 95 (7.4, bis 30 % zu hoch LL40M, 30°-Tangente);
%   Has18 S. 29 (3.2 FE-Modell, 20 Stellungen); Kas22 S. 6 bis 9 (4.1, 4.2 Hill), S. 7 (Biegespannung lokal),
%   S. 12 (5.1.1 VDI 2736 etwa 20 % zu hoch VDI2), S. 14 (5.1.2 bis 25 % zu niedrig VDI1), S. 17 (Summary).
% Absatz 4 (Ellipse): Frue12 S. 93 (5.4, Abbildung 41), S. 94 bis 98 (5.4.1, Gl. 89 bis 100), S. 99 bis 101
%   (5.4.2, 19 Verzahnungen, γ und a/b, Abbildung 46, 11,9 %, 14,5 %, 14,8 %, Spline 3 %), S. 102 (Tragfähigkeit,
%   Volumen), S. 103 (Exkurs 12 bis 30 %); Kas22 S. 3 bis 4 (3.1, Gl. 1 bis 8, d_f,e = d_f + m_n e_f), S. 14
%   (e_f = −0,2 und −0,4, kleinstes h*_fP); Lan21 S. 7 bis 8 (3.1, 3.2 D1, D2), S. 10 bis 11 (Tabellen 4 und 5,
%   etwa 15 %), S. 16 (Table 10, < 5 %), S. 17 (10 bis 15 % Last); San10 S. iii (Abstract 7,6/22,4/24,3 %),
%   S. iv (mehr als 30-fach), S. 81 (33-fach).
% Absatz 5 (Bézier): DE102013004861B3 S. 2 (Zusammenfassung, [0007]), S. 3 ([0011] bis [0016]), S. 4 ([0018] bis
%   [0020]), S. 6 bis 7 ([0041] bis [0044], Fig. 2), S. 8 bis 9 ([0054] bis [0062], rund 35 %); RotOpf17 S. 720
%   (bionisch bis 30 %, Patent 2009), S. 722 (Bézier-Faktor, Kathete des ersten Zugdreiecks), S. 724 (Fig. 6,
%   Optimum 0,57, 24 %, 0,46 bis 0,81), S. 725 (1,25 auf 1,55), S. 728 bis 729 (Fig. 10, 12 bis 30 %), S. 732
%   (Summary); FVA241-XII S. 32 bis 35 (2.4.2 Roth), S. 42 bis 43 (3.2.2, Abbildung 3.2, Be 0,5 bis 0,6 etwa 25 %,
%   Be = 0,57), S. 93 (5.1.3, Abbildung 5.17, etwa 20 %), S. 94 bis 95 (Tabelle 5.7, ≤ 3 %, 36 bis 37°);
%   Don20 S. 1 (Abstract), S. 4 (Fig. 2 Bézier mit fünf Punkten), S. 12 (18,19 %, 20,06 %, 22,02 %), S. 15.
% Absatz 6 (bionisch): DE102008045318B3 S. 2 ([0005], [0006] bis 30 %), S. 3 ([0018] Koordinaten, [0022] d_r als
%   Mittel aus d_N und d_FR), S. 4 ([0023] γ, [0024] Übergangspunkt, [0025] Gl. 1), S. 5 ([0040] γ < 65°, ≤ 45°,
%   Gl. 15), S. 6 ([0049] γ und r frei), S. 7 ([0050] Radius 0,1 bis 0,6 bzw. 0,3 bis 0,4 S_L, [0051] FEM bis 30 %),
%   S. 9 (Anspruch 5: 45° − 180°/z − α_r), S. 10 (Anspruch 12); Kas22 S. 5 bis 6 (3.3, Gl. 14 bis 17), S. 14
%   (5.1.2 γ_b ≈ 30°, kleines b_f), S. 16 (Figure 16: 7,3 %, 4,84 %, 2,63 %), S. 16 bis 17 (5.2 Verschleißkerbe),
%   S. 17 (verstärkt/unverstärkt); Kas23 S. 2596 (CAO, Gl. 1), S. 2597 (δ_per = 0,025 m_n Gl. 3, Figure 7 und 8,
%   Kontakt des Gegenrads), S. 2598 (18,7 %), S. 2599 (Reibwert, Vergleichsspannung, 17 bis 24 %, MOTT = Patent),
%   S. 2600 (Summary).
%   Abgleich mit fillet.py: Ellipse ab d_Ff zwischen P_Ff und P_f (Frühe Abbildung 41); bionisch und Bézier ab
%   d_r = Mittel aus d_N und d_FR (DE 10 2008 045 318 B3 Anspruch 12 und [0022]); Lückenkoordinaten [0018];
%   Nomenklatur der Arbeit: φ für den Patentwinkel γ, f und l für die Bézierkontur.
% 2.5.3 Absatz 1: Sto18 S. 307 (abschätzender Charakter, Machbarkeitsstudien), S. 308 (6.1.1 Oberbach, Kurven),
%   S. 309 (Erweiterungen, formelle Struktur), S. 310 bis 313 (6.1.2 lokale Bauteilfestigkeit, Bild 6.1 bis 6.3),
%   S. 311 (Auslastungsgrad = Kehrwert des Sicherheitsfaktors), S. 313 (Gl. 6.1), S. 322 (6.1.3 FKM-Analogie).
% Absatz 2: Sto18 S. 322 (Gl. 6.2 = eq:q_ftau), S. 323 (f_τ 0,5/1,0/0,58, Tabelle 6.1 mit Fußnote PA+GF),
%   S. 324 (Mehraxialität, spröde), S. 325 (Gl. 6.3 = eq:h_hencky, Tabelle 6.2, 2/3 an der Oberfläche), S. 326
%   (Gl. 6.4 = eq:eps_zul), S. 344 bis 346 (6.2.2 Nomogramm Bild 6.17, K = Streckgrenze, Bild 6.18 Faktoren),
%   S. 345 bis 346 (faserverstärkt pauschal, isotrope FE, Kurve in Faserrichtung), S. 347 bis 350 (6.2.2.1
%   Temperatur, Bild 6.21), S. 355 (Gl. 6.15), S. 356 (Vergleichsdehnung abgelesen), S. 357 (Gl. 6.16, 6.17 =
%   eq:druckkorrektur, 30 % höhere Druckfestigkeit), S. 358 (Sicherheitsfaktoren implizit).
% Absatz 3: Sto18 S. 328 (6.1.4 Grenzspannungsamplitude N_G = 10^7, weiterer Abfall), S. 328 bis 329 (Goodman
%   gegen Gerber Gl. 6.7, Normierung), S. 360 (Gl. 6.18, Bild 6.24), S. 361 (HCF > 10^4, statischer Nachweis als
%   Voraussetzung), S. 362 (6.3.1 linear-elastisch, 6.3.2, Bild 6.26 A_W), S. 363 (Gl. 6.19, feste Werte,
%   Wechselfestigkeitsfaktor, Stähle 0,4 bis 0,45), S. 364 (Bild 6.27 POM, PA66-GF35), S. 364 bis 365 (6.3.3,
%   Bild 6.28 Haigh, M = σ_W/σ_Y), S. 366 bis 367 (Überlastungsfälle Bild 6.29, Gl. 6.20, Wahl durch Anwender),
%   S. 367 bis 368 (Gl. 6.21 = eq:k_m_ueberlastung, M = σ_W/σ_Y).
% Absatz 4: Sto18 S. 369 bis 370 (6.3.4, Bild 6.31, Basquin, Verschiebung), S. 371 (Gl. 6.22, Tabelle 6.3 k 15,65
%   bis 20,44), S. 372 (11,9 bis 25,0; Grenze 10^4; Gl. 6.23), S. 373 bis 374 (6.3.5, Gl. 6.24 bis 6.26, q ohne
%   Druckkorrektur), S. 375 (Faktor 2 bis 3, Stützzahlen, S-Life).
% Absatz 5: Sto18 S. 343 bis 344 (6.2.1 isotrop, Prüfstab in Fließrichtung, anisotrop unzulässig, Reduktions-
%   faktoren unzulässig), S. 344 bis 346 (Minderung quer über Abminderungsfaktor), S. 363 (A_W bei verstärkt nicht
%   h-abhängig), S. 235 (3.3.4 umfassender Nachweis nur mit Software, Kapitel 6 als Alternative), S. 375.
% Absatz 6: Sto18 S. 219 (3.3 Richtung und Orientierungsgrad), S. 221 (3.3.1 Versagensmechanismen), S. 226 bis 227
%   (3.3.2 Platten 0°/90°/30° oder 45°), S. 229 bis 234 (3.3.3 Kriterien), S. 231 bis 232 (Tsai-Hill 2D Gl. 3.34
%   und 3.35), S. 233 (Gl. 3.36 = eq:tsai_hill, transversal isotrop vier Parameter, a_TH = Wurzel), S. 236 bis 237
%   (3.3.4.2 Gl. 3.40), S. 238 bis 239 (3.3.4.3 Interpolation), S. 241 bis 243 (3.3.4.5 lokales Wöhlerkurven-
%   Konzept, Converse, komponentenweise, Tsai-Hill-Überlagerung, S-Life), S. 242 (linear-elastisch im HCF);
%   Vertauschung L/N: eq:hill in kap3_7_auswertung.tex führt L mit σ23 und N mit σ12 (Abaqus), das Buch L mit τ12.
% Bilder fig:nomogramm_qs und fig:haigh_qs_80C wörtlich aus kap2_stand.tex übernommen.
% BA-0486-Prüfung: Gl. (2.1) → Nie19 S. 58; (2.3) → VDI Gl. 10 (bestätigt); Tabelle 2.2 → S_Fmin 2,0/1,6 (VDI S. 14,
%   S_Hmin hier nicht übernommen); (2.10) → Sto18 Gl. 6.3; (2.11) → Gl. 6.15 (steht in 3.7); (2.12) → Gl. 6.2;
%   (2.13) bis (2.16) → Gl. 6.16/6.17 (bestätigt, f_σ 1,3); Tabelle 2.5 → Tabelle 6.1 S. 323 (Werte bestätigt);
%   (2.17) Basquin → S. 370 (im Buch nur als Aussage); (2.18) bis (2.20) → Gl. 6.18/6.22 (in 3.7); (2.21) → Gl. 6.21
%   Fall F2 (bestätigt); (2.22) → Gl. 6.19/6.20 (bestätigt); BA-Aussage "Oberbach erweitert um Mittelspannung" →
%   Sto18 S. 309.
% Bildwunsch: kap2/quelle_fruehe_ellipse aus 914751_Fruehe_2012.pdf, Druckseite 93 (PDF 105), Abbildung 41
%   "Zahnfußausrundung mit tangentenstetigen Ellipsen"; Ausschnitt die Skizze mit P_Ff, M, a, b, γ, P_f und den
%   Beschriftungen "tangentenstetiger Übergang".
% Bildwunsch: kap2/quelle_patent_bezier aus DE102013004861B3.pdf, Fig. 2 auf den Zeichnungsblättern (Seite 12/15,
%   Blatt prüfen): symmetrische Zahnlücke mit P0 bis P3, Q0 bis Q2, t1, t2, S, d_r; nur Fig. 2 ausschneiden.
% Bildwunsch: kap2/quelle_patent_bionisch aus DE102008045318B3.pdf, Fig. 2 auf den Zeichnungsblättern (Seite 11/12,
%   Blatt prüfen): Zahnfußform mit Tangensfunktion 5 und 5', Kreisbahn 6, ÜP, FP, γ, d_r; nur Fig. 2 ausschneiden.
% Bibwunsch: NieWin03, Fuer13, San10, FVA30-XI, FVA30-XII, FVA892-I, FVA484-V, FVA241-XII in bib_kap2_5.bib;
%   vorhanden in litbib.bib: VDI2736-2, ISO6336-3, Nie19, Lin10, Has18, Frue12, Sto18, Kas22, Kas23, Lan21, Don20,
%   RotOpf17, DE102008045318B3, DE102013004861B3.
% Nomenklatur: sigma_F, MPa, örtliche Zahnfußspannung (VDI), VDI2736-2 S. 11 (fehlt)
% Nomenklatur: sigma_FP, MPa, zulässige Zahnfußspannung, VDI2736-2 S. 13 (fehlt)
% Nomenklatur: sigma_FG, MPa, Zahnfußgrenzfestigkeit, VDI2736-2 S. 13 (fehlt)
% Nomenklatur: sigma_FlimN, MPa, Zeitschwellfestigkeit als Nennspannung, VDI2736-2 S. 14 (fehlt)
% Nomenklatur: K_F, -, Faktor für Zahnfußbeanspruchung, VDI2736-2 S. 11 (fehlt); K_A, -, Anwendungsfaktor, S. 11 (fehlt)
% Nomenklatur: Y_Fa, -, Formfaktor für Kraftangriff am Zahnkopf, VDI2736-2 S. 11 (fehlt)
% Nomenklatur: Y_Sa, -, Spannungskorrekturfaktor für Kraftangriff am Zahnkopf, VDI2736-2 S. 12 (fehlt)
% Nomenklatur: Y_epsilon, -, Überdeckungsfaktor, VDI2736-2 S. 13 (fehlt); Y_beta, -, Schrägenfaktor, S. 13 (fehlt)
% Nomenklatur: Y_St, -, Spannungskorrekturfaktor der Prüfräder (etwa 2,0), VDI2736-2 S. 13 (fehlt)
% Nomenklatur: S_Fmin, -, geforderte Mindestsicherheit Zahnfuß, VDI2736-2 S. 14 (fehlt)
% Nomenklatur: F_t, N, Nennumfangskraft am Teilzylinder, VDI2736-2 S. 11 (fehlt)
% Nomenklatur: N_L, -, Lastwechselzahl, VDI2736-2 S. 14 (fehlt; N geforderte Lastspielzahl nach Sto18 ist gelistet)
% Nomenklatur: vartheta_Fuß, °C, Zahnfußtemperatur, VDI2736-2 S. 8 (fehlt)
% Nomenklatur: z_n, -, Ersatzzähnezahl, VDI2736-2 S. 11 (fehlt)
% Nomenklatur: h_fP, mm, Fußhöhe des Bezugsprofils; rho_fP, mm, Fußrundungsradius des Bezugsprofils; VDI2736-2 S. 12
%   (fehlen; h_aP0 und rho_aP0 des Werkzeugs sind gelistet)
% Nomenklatur: Y_F, -, Formfaktor Methode B, ISO6336-3 S. 17 (fehlt); Y_S, -, Spannungskorrekturfaktor Methode B, S. 24 (fehlt)
% Nomenklatur: s_Fn, mm, Zahnfußdickensehne; h_Fe, mm, Biegehebelarm; rho_F, mm, Fußrundungsradius an der 30°-Tangente;
%   ISO6336-3 S. 19 (fehlen)
% Nomenklatur: alpha_Fen, °, Kraftangriffswinkel, ISO6336-3 S. 17 (fehlt); theta, °, Tangentenwinkel, S. 19 (fehlt)
% Nomenklatur: f_epsilon, -, Faktor der Lastaufteilung in Gl. (9), ISO6336-3 S. 17 (fehlt)
% Nomenklatur: q_s, -, Kerbparameter s_Fn/(2 rho_F), ISO6336-3 S. 25 (fehlt)
% Nomenklatur: L, -, Verhältnis s_Fn/h_Fe, ISO6336-3 S. 24 (fehlt; kollidiert mit der Hill-Konstante L, Nutzer entscheidet)
% Nomenklatur: epsilon_alpha,w, -, Profilüberdeckung unter Last, Has18 S. 31 (fehlt; epsilon_alpha gelistet)
% Nomenklatur: P_Ff, P_f, Fußformkreispunkt und Fußpunkt der Ellipse, Frue12 S. 93 (fehlen; Punkte, keine Größen)
% Nomenklatur: alpha_r, °, Eingriffswinkel am relevanten Durchmesser, DE102008045318B3 S. 5 (fehlt)
% Nomenklatur: P_0 bis P_3, S, Haupt-, Kontroll- und Schnittpunkt der Bézierkurve, DE102013004861B3 S. 6 (nur Bildunterschrift)
% Nomenklatur: sigma_H, MPa, hydrostatische Spannung, Sto18 S. 325 (fehlt)
% Nomenklatur: epsilon_max, epsilon_min, %, obere und untere zulässige Dehnung aus dem Abminderungsfaktorbereich, Sto18 S. 326 (fehlen)
% Nomenklatur: sigma_W' (gestrichen), MPa, mittelspannungskorrigierte Wechselfestigkeit, Sto18 S. 367 (fehlt)
% Nomenklatur: sigma_min, sigma_max, MPa, Minimal- und Maximalspannung des Lastspiels, Sto18 S. 368 (fehlen)
% Nomenklatur: R, -, Spannungsverhältnis, Sto18 S. 365 (fehlt, wie in 3.7 gemeldet)
% Nomenklatur: X, Y, Z, MPa, Zugfestigkeiten in den Orthotropierichtungen; S_ij, MPa, Schubfestigkeiten; Sto18 S. 233 (fehlen)
% Nomenklatur: gelistet sind b, m_n, z, alpha_n, epsilon_alpha, a/b, gamma, varphi, r, S_L, d_r, d_N, d_FR, f, l, q, f_tau,
%   h, epsilon_zul, K, A, A_W, sigma_stat, sigma_W, sigma_Y, M, K_M, K_B, N_G, N, k, sigma_m, sigma_a, sigma_1, sigma_3,
%   sigma_V,Mises, f_sigma, sigma_ij, F, G, H, L, M, N; sigma_2 fehlt (wie in 3.7 gemeldet)
% Nomenklatur: Abkürzungen ACORA (Has18 S. 35), CAO (Kas23 S. 2592), LDP (San10 S. iii), RIKOR, STIRAK, OeZS
%   (FVA30-XI S. 11), FKM (Sto18 S. 322) fehlen in abkuerzungen.
% Offen: BA-Quellen Hac66, Nie60, Ohl58, Ehr11, DIN11, Kar10, Tho12, Has16, Rei22, Nie03, ISO19a liegen nicht vor:
%   Hachmann/Strickle nur als Nennung über VDI 2736 Anhang B (S. 30); Zahnverlustgrad nach Ohlendorf (BA Gl. 2.4, 2.5),
%   relative Einschaltdauer nach Ehrenstein/DIN EN 60034 und BA Gl. 2.6 bis 2.8 samt Tabellen 2.3 und 2.4 gestrichen
%   (Zahntemperatur nur ein Satz); Karimpour und Thoma durch Has18 S. 31 ersetzt (Thoma als Nennung in Has18 S. 35);
%   Hasl et al. 2016 durch Has18 ersetzt; Reitschuster (PEEK) entfällt; ISO-6336-Lebensdauerfaktoren Y_NT, Z_NT und
%   die Flankenformel (BA Gl. 2.2) gestrichen (Thema Fuß); BA-Satz zur Ausfallwahrscheinlichkeit auf Has18 S. 22
%   gestützt; JIS B 1759, BS 6168 (Sta20, Wal89) gestrichen.
% Offen: Ren20 (FKM) nur als Nennung über Sto18 S. 322 und 363; Obe81 und Vor10 nur über Sto18 S. 308 und 325
%   genannt; Ehr07 ohne Textebene nicht verwendet (Wechselfestigkeitsfaktor über Sto18 S. 363).
% Offen: Paper Frühe "Minimierung von Zahnfußspannungen": FrueheMinimierung.txt ist mit Fruehe2012.txt identisch
%   (cmp), daher nur die Dissertation zitiert.
% Offen: KassemAnalytical (Laminattheorie, Anbindung an VDI 2736): Kopfzeile ohne Jahr und Journal, PDF-Metadaten
%   "Manuskriptvorlage VDI Wissensforum", Dezember 2024; nicht zitiert, Bib-Eintrag unvollständig.
% Offen: FVA 241 X (Normberechnung asymmetrische Verzahnung) trägt zu den Fußgeometrien nichts bei, nicht zitiert.
% Offen: FVA 30 XI: Ziffern des Titelblatts (Heft, Jahr) fehlen in der Textebene (Schriftcodierung); Heft 1586 und
%   2023 aus dem Briefing, am Deckblatt prüfen; Autoren = Sachbearbeiter Schweigert und Hofer (decodiert).
% Offen: Fuer13 liegt als "Entwurf: Stand 01.03.2013" vor (Einreichdatum "XYZ" im Titelblatt); Seiten dieser Fassung.
% Offen: RotOpf17-Seiten als 718 + PDF-Seite angenommen; am Original prüfen.
% Offen: Delrin 100 NC010 (FVA892-I S. 28) ist ein POM, im Text als Handelsname belassen.
% Offen: Y_epsilon und Y_beta mit griechischen Indizes in der Schreibweise der Richtlinie (nicht aufrecht).
% Offen: Dreier-Hinweise des Textchecks sind Aufzählungen aus den Quellen (Faktoren, Kriterien, Versagensmodi).
